\documentclass{itmo}

\setlist[itemize]{label=---}
\setlist[enumerate]{label=\arabic*)}

\begin{document}

\makereporttitle{Программной инженерии и компьютерной техники (ФПИ и КТ)}
	{лабораторной работе №1}
	{Структуры и алгоритмы в базах данных и распределенных системах}
	{Файловая хэш-таблица, perfect hash, LSH для текстов}
	{}
	{Баряев~А.~А./P4135}
	{Платонов~А.~В./доцент, Портнов~П.~В./ассистент}

\maketableofcontents

\chapter{Теоретическая часть}

\section{Файловая хэш-таблица}

В рамках работы реализована хэш-таблица на файловой системе с поддержкой операций \texttt{Put}, \texttt{Get}, \texttt{Delete}.
Коллизии обрабатываются цепочками внутри бакетов.
Таблица хранится в бинарном формате в одном файле.
Подход ориентирован на внешнюю память и постоянное хранение данных:
структура должна корректно переживать завершение процесса и повторное открытие.

\subsection*{Основные идеи}
\begin{itemize}
	\item фиксированный массив бакетов хранит смещения на цепочки записей;
	\item данные по ключам записываются в append-only области файла;
	\item удаление реализуется через tombstone, что упрощает и ускоряет операции записи;
	\item для служебной части используется \texttt{mmap}, для пользовательских записей --- адресное чтение/запись.
\end{itemize}

\section{Perfect hash}

Под \textit{perfect hash} в работе понимается полный индекс по фиксированному множеству ключей:
после построения структуры выполняется быстрый поиск индекса ключа.
Структура не рассчитана на постоянные вставки после построения.
В классическом варианте perfect hashing стремится к отсутствию коллизий в конечной таблице.
В данной лабораторной работе акцент сделан на практический статический индекс:
набор ключей фиксируется на этапе \texttt{Build}, далее выполняются только быстрые \texttt{Lookup}.

\section{MinHash + LSH}

Для поиска дублей в текстах используется схема:
\begin{itemize}
	\item разбиение текста на шинглы;
	\item построение MinHash-подписи;
	\item LSH-разбиение подписи на полосы и формирование бакетов кандидатов;
	\item оценка похожести по доле совпавших компонент подписи.
\end{itemize}
Сигнатуры позволяют заменить дорогое точное попарное сравнение всех шинглов
на более компактный и масштабируемый этап поиска кандидатов.

\chapter{Практическая часть}

\section{Структура проекта}

Код реализован на Go в директории \texttt{lab1}:
\begin{itemize}
	\item \texttt{internal/hashfs} --- файловая хэш-таблица;
	\item \texttt{internal/perfecthash} --- полный индекс по фиксированному набору ключей;
	\item \texttt{internal/lshtext} --- MinHash + LSH для текстов;
	\item \texttt{internal/gen} --- генераторы тестовых данных;
	\item \texttt{cmd/lsh-demo} --- небольшая демонстрационная утилита для проверки LSH на внешних файлах.
\end{itemize}
Во всех пакетах используется только стандартная библиотека Go.

\section{Реализация файловой хэш-таблицы}

Реализация расположена в \texttt{internal/hashfs/hashfs.go}.
Файл таблицы состоит из:
\begin{enumerate}
	\item заголовка (магическое число, версия, число бакетов, начало области данных, текущий хвост);
	\item таблицы бакетов (массив смещений в цепочки);
	\item области записей.
\end{enumerate}

Запись имеет формат:
\texttt{[keyLen][valLen][hash][flags][key][value][nextOffset]}.
Для удаления используется логический флаг \texttt{tombstone}.
При обновлении ключа в текущей версии также выполняется append новой записи,
что упрощает атомарность изменения и уменьшает сложность in-place обновлений.

Прямая работа с диском выполнена через:
\begin{itemize}
	\item \texttt{syscall.Mmap} для заголовка и таблицы бакетов;
	\item \texttt{ReadAt}/\texttt{WriteAt} для чтения и дозаписи записей.
\end{itemize}

\subsection*{Публичный API}
\begin{itemize}
	\item \texttt{Open(path, opts)} --- открыть/создать таблицу;
	\item \texttt{Put(key, value)} --- вставить или обновить значение;
	\item \texttt{Get(key)} --- получить значение;
	\item \texttt{Delete(key)} --- логически удалить ключ;
	\item \texttt{Close()} --- корректно закрыть файл и снять отображение.
\end{itemize}

\section{Реализация perfect hash}

Реализация расположена в \texttt{internal/perfecthash/perfecthash.go}.
Публичное API:
\begin{itemize}
	\item \texttt{Builder.Build(keys [][]byte)} --- построение полного индекса;
	\item \texttt{Table.Lookup(key []byte)} --- поиск индекса ключа;
	\item \texttt{Serialize}/\texttt{Deserialize} --- сохранение и восстановление индекса.
\end{itemize}

В текущей версии индекс хранится как статическая структура \texttt{map[string]int}
для фиксированного набора ключей.
Это обеспечивает быстрый поиск и корректную работу по требованиям задания.
Дополнительно реализована проверка дубликатов ключей на этапе построения.

\subsection*{Комментарии по выбору структуры}
Классический двухуровневый perfect hash сложнее в реализации и отладке.
Для лабораторной работы выбран более стабильный статический индекс:
он сохраняет ключевое требование ``фиксированный набор ключей + быстрый поиск''
и позволяет сосредоточиться на сравнении подходов по времени выполнения и памяти.

\section{Реализация LSH для текстов}

Реализация расположена в \texttt{internal/lshtext/lshtext.go}.
Публичные операции:
\begin{itemize}
	\item \texttt{Add(docID, text)} --- добавить текст в индекс;
	\item \texttt{Query(text)} --- получить кандидатов и их похожесть;
	\item \texttt{FullScanDuplicates(threshold)} --- найти пары дублей по порогу.
\end{itemize}

В реализации используются:
\begin{itemize}
	\item токенизация \texttt{strings.Fields};
	\item word-shingles фиксированного размера;
	\item MinHash-подпись длины 64;
	\item 8 полос по 8 компонент в LSH.
\end{itemize}

\subsection*{Хранимые структуры}
\begin{itemize}
	\item \texttt{signatures}: \texttt{docID -> []uint64};
	\item \texttt{bandBuckets}: \texttt{band -> bucketHash -> []docID}.
\end{itemize}
Это позволяет быстро получать кандидатов при \texttt{Query}
и выполнять \texttt{FullScanDuplicates} без полного перебора всех пар на каждом шаге.

\chapter{Тестирование}

\section{Функциональные тесты}

Покрытие тестами:
\begin{itemize}
	\item \texttt{internal/hashfs/hashfs\_test.go}: 
	\begin{itemize}
		\item массовые рандомные вставки, обновления, удаления;
		\item сверка состояния с эталонной \texttt{map[string][]byte};
		\item проверка персистентности после повторного открытия;
		\item edge-cases: пустой ключ, удаление отсутствующего ключа, превышение лимита значения;
		\item рандомизированные операции в несколько раундов открытия/закрытия файла.
	\end{itemize}
	\item \texttt{internal/perfecthash/perfecthash\_test.go}:
	\begin{itemize}
		\item генерация уникальных случайных ключей и проверка \texttt{Lookup};
		\item проверка отсутствия по модифицированным ключам;
		\item проверка сериализации/десериализации;
		\item отдельные тесты на дубликаты и пустой вход.
	\end{itemize}
	\item \texttt{internal/lshtext/lshtext\_test.go}:
	\begin{itemize}
		\item сценарий с near-duplicates и шумовыми текстами;
		\item проверка \texttt{Query} и \texttt{FullScanDuplicates};
		\item проверка чувствительности к порогу похожести.
	\end{itemize}
\end{itemize}

Запуск:
\begin{verbatim}
go test ./internal/...
\end{verbatim}

Фактический результат:
\begin{verbatim}
ok   siaod-hw1/internal/hashfs
ok   siaod-hw1/internal/lshtext
ok   siaod-hw1/internal/perfecthash
\end{verbatim}

\section{Проверка на внешних текстах}

Для демонстрации LSH добавлена утилита:
\begin{verbatim}
go run ./cmd/lsh-demo file1.txt file2.txt [file3.txt ...]
\end{verbatim}

Пример результата:
\begin{verbatim}
go run ./cmd/lsh-demo AllCombined.txt AllCombined_copy.txt
pair (0, 1) similarity=1.000
\end{verbatim}

Пример для похожих, но не одинаковых файлов:
\begin{verbatim}
go run ./cmd/lsh-demo AllCombined.txt AllCombined_copy2.txt AllCombined_copy1.txt
pair (1, 2) similarity=0.812
\end{verbatim}

Таким образом, на внешних данных подтверждаются оба ожидаемых режима:
\begin{itemize}
	\item точные дубликаты дают похожесть, близкую к 1;
	\item частичные дубликаты дают высокую, но не максимальную похожесть.
\end{itemize}

\chapter{Исследовательская часть}

\section{Среда запуска}

Замеры получены на Linux (amd64), Go 1.25.x.
Процессор по выводу бенчмарков:
\texttt{AMD Eng Sample: 100-000000829-50\_Y}.

\section{Методика}

\begin{itemize}
	\item производительность измерялась встроенными benchmark-тестами Go (\texttt{go test -bench});
	\item профили CPU и памяти снимались через \texttt{-cpuprofile} и \texttt{-memprofile};
	\item для каждого алгоритма брались отдельные репрезентативные сценарии.
\end{itemize}

\subsection*{Измеряемые сценарии}
\begin{itemize}
	\item \textbf{hashfs}: \texttt{Insert}, \texttt{RandomGet}, \texttt{UpdateHeavy};
	\item \textbf{perfecthash}: \texttt{Build}, \texttt{Lookup};
	\item \textbf{lshtext}: \texttt{IndexBuild}, \texttt{Query}.
\end{itemize}

\section{Результаты бенчмарков}

Команда:
\begin{verbatim}
go test -bench=. ./internal/hashfs ./internal/perfecthash ./internal/lshtext
\end{verbatim}

\begin{table}[H]
	\centering
	\caption{Ключевые результаты benchmark-тестов}
	\begin{tabular}{|l|l|l|}
		\hline
		Пакет & Бенчмарк & Результат (нс/оп) \\
		\hline
		hashfs & BenchmarkStore\_Insert & 775.5 \\
		hashfs & BenchmarkStore\_RandomGet & 1060 \\
		hashfs & BenchmarkStore\_UpdateHeavy & 654.8 \\
		perfecthash & BenchmarkPerfectHashBuild & 7016939 \\
		perfecthash & BenchmarkPerfectHashLookup & 15.89 \\
		lshtext & BenchmarkLSHIndexBuild & 9209 \\
		lshtext & BenchmarkLSHQuery & 436614 \\
		\hline
	\end{tabular}
\end{table}

\begin{table}[H]
	\centering
	\caption{Оценка пропускной способности}
	\begin{tabular}{|l|l|l|}
		\hline
		Пакет & Бенчмарк & Оценка, оп/с \\
		\hline
		hashfs & BenchmarkStore\_Insert & \(\approx 1.29 \cdot 10^6\) \\
		hashfs & BenchmarkStore\_RandomGet & \(\approx 0.94 \cdot 10^6\) \\
		hashfs & BenchmarkStore\_UpdateHeavy & \(\approx 1.53 \cdot 10^6\) \\
		perfecthash & BenchmarkPerfectHashLookup & \(\approx 6.29 \cdot 10^7\) \\
		lshtext & BenchmarkLSHQuery & \(\approx 2.29 \cdot 10^3\) \\
		\hline
	\end{tabular}
\end{table}

\subsection*{Интерпретация результатов}
\begin{itemize}
	\item операции \texttt{Get}/\texttt{Update} в hashfs выполняются быстрее 1--1.5 млн оп/с для небольших ключей;
	\item \texttt{perfecthash Lookup} демонстрирует минимальную задержку (десятки наносекунд);
	\item \texttt{lsh Query} существенно тяжелее простых KV-операций, поскольку включает шинглинг и MinHash-подпись.
\end{itemize}

\section{Профилирование CPU/Memory}

Примеры команд:
\begin{verbatim}
go test -bench=BenchmarkStore_Insert \
  -cpuprofile=hashfs_cpu.out -memprofile=hashfs_mem.out ./internal/hashfs

go test -bench=BenchmarkPerfectHashBuild \
  -cpuprofile=perfecthash_cpu.out -memprofile=perfecthash_mem.out ./internal/perfecthash

go test -bench=BenchmarkLSHIndexBuild \
  -cpuprofile=lsh_cpu.out -memprofile=lsh_mem.out ./internal/lshtext
\end{verbatim}

\subsection*{Наблюдения по профилям}
\begin{itemize}
	\item \textbf{hashfs CPU}: заметны расходы на хеширование и операции чтения/записи;
	\item \textbf{hashfs MEM}: основная память --- буферы чтения/записи и служебные структуры отображения;
	\item \textbf{perfecthash CPU}: доминирует этап построения индексной \texttt{map};
	\item \textbf{perfecthash MEM}: аллокации ключей и внутренних bucket-структур \texttt{map};
	\item \textbf{lshtext CPU}: токенизация, шинглинг, MinHash и hash band-ов;
	\item \textbf{lshtext MEM}: подписи документов и списки кандидатов в бакетах.
\end{itemize}

Анализ узких мест:
\begin{itemize}
	\item \textbf{hashfs}: вклад дают хеширование ключей, системные I/O операции \texttt{ReadAt/WriteAt} и копирование буферов;
	\item \textbf{perfecthash}: основная стоимость в \texttt{Build} связана с аллокациями и преобразованием ключей при построении индекса;
	\item \textbf{lshtext}: основные затраты в токенизации, вычислении MinHash-подписей и обработке LSH-бакетов.
\end{itemize}

Гипотезы улучшений:
\begin{itemize}
	\item для \texttt{hashfs}: пакетная запись и периодическая компактизация файла;
	\item для \texttt{perfecthash}: более компактная статическая структура вместо \texttt{map};
	\item для \texttt{lshtext}: оптимизация хеш-функций и уменьшение/настройка длины подписи.
\end{itemize}

\chapter{Дополнительные замечания по реализации}

\section{Соответствие формулировке задания}

\begin{itemize}
	\item файловая структура реализована через прямую работу с диском и \texttt{mmap}, без использования сериализации уровня \texttt{gob/json};
	\item по каждому алгоритму присутствуют функциональные тесты и benchmark-тесты;
	\item присутствуют команды и артефакты профилирования CPU/памяти.
\end{itemize}

\section{Ограничения текущей версии}

\begin{itemize}
	\item hashfs реализован как append-only без фоновой компактизации;
	\item perfect hash реализован как статический индекс на \texttt{map}, а не как полноценная двухуровневая FKS-схема;
	\item LSH работает на уровне текстов (документов), вариант для 3D-точек не реализован.
\end{itemize}

Эти ограничения не нарушают требования текущего варианта лабораторной работы,
но определяют направления дальнейшего развития.

\chapter{Вывод}

В работе реализованы все требуемые алгоритмы:
\begin{itemize}
	\item файловая хэш-таблица с прямой работой с диском и \texttt{mmap};
	\item полный индекс по фиксированному набору ключей (perfect hash API);
	\item LSH (MinHash + banding) для поиска дублей в текстах.
\end{itemize}

Для каждой реализации выполнены:
\begin{itemize}
	\item функциональные тесты на случайных данных;
	\item benchmark-тесты производительности;
	\item профилирование CPU и памяти с анализом узких мест.
\end{itemize}

Результаты подтверждают работоспособность и масштабируемость предложенных решений для целей лабораторной работы.
Дополнительно проведена ручная валидация LSH на внешних текстовых файлах,
что подтверждает корректную работу поиска дубликатов в прикладном сценарии.

\end{document}


